% Generated draft: requires mathematical and source review.
\documentclass[11pt]{article}
\usepackage[margin=25mm]{geometry}
\usepackage{fontspec}
\setmainfont{Latin Modern Roman}
\usepackage{amsmath,amssymb,mathtools,unicode-math}
\setmathfont{Latin Modern Math}
\usepackage{enumitem,needspace,xurl,hyperref}
\setlength{\parindent}{0pt}
\setlength{\parskip}{5pt}
\setlength{\emergencystretch}{3em}
\newcommand{\N}{\mathbb{N}}
\newcommand{\Z}{\mathbb{Z}}
\newcommand{\Q}{\mathbb{Q}}
\newcommand{\R}{\mathbb{R}}
\newcommand{\C}{\mathbb{C}}
\newcommand{\F}{\mathbb{F}}
\newcommand{\A}{\mathbb{A}}
\newcommand{\PP}{\mathbb{P}}
\newcommand{\E}{\mathbb{E}}
\newcommand{\eps}{\varepsilon}
\newcommand{\dd}{\,\mathrm{d}}
\begin{document}
\section*{UnsolvedMath ID: 30005301}
\textbf{Source problem code:} catalogue entry
\subsection{Definitions and mathematical statement}
For terminology only, let \(k\) be an algebraically closed field and let \(Z\subset\PP^n_k\) be a finite reduced set of points. The set \(Z\) is in \emph{linear general position} if every subset of \(r\leq n+1\) points spans a projective subspace of dimension \(r-1\).

Under the usual spatial convention, a finite reduced set \(Z\subset\PP^3_k\) is \emph{geproci} (general projection complete intersection) if there are positive integers \(a,b\) and a nonempty Zariski open set of projection centers \(p\in\PP^3_k\setminus Z\) such that projection from each such \(p\) is injective on \(Z\), and its image in \(\PP^2_k\) is the scheme-theoretic complete intersection of two plane curves of degrees \(a,b\) having no common component. In particular, \(|Z|=ab\). This definition supplies a conventional meaning of the terminology; the supplied excerpt does not establish that this ambient dimension or these field hypotheses are the intended ones.

The source question is: ``Is it true that such point configurations are never geproci?'' If \(\mathcal C\) denotes the class referred to by ``such point configurations,'' its logical form is
\[
\text{Is it true that }\quad
\forall Z\in\mathcal C,\qquad Z\text{ is not geproci}?
\]
The antecedent defining \(\mathcal C\) is missing. The supplied title and dated assessment associate the question with configurations in linear general position and with nontrivial examples having more than four points. They do not recover the ambient projective space, the field and characteristic, the exact meaning of ``nontrivial,'' or all restrictions on the configurations or complete-intersection degrees. Thus they do not determine a precise class \(\mathcal C\), and no unconditional quantified nonexistence conjecture can be recovered from this record.

The supplied registry marks the statement as \texttt{unrecoverable}. Its verification note describes it as a narrower duplicate of the geproci classification questions retained in \texttt{OWR-11695864-001}, and explicitly declines to promote a separate clean statement. That other record is not supplied here. The reciprocal-plane and Orlik--Terao passage following the question concerns a different topic and does not supply the missing antecedent.
\subsection{Short English statement}
The record asks whether a context-dependent class of point configurations can ever be geproci. Its title and supplied assessment suggest linear general position and nontrivial examples with more than four points, but the defining context is missing. The conventional terminology is explained above; the intended conjecture remains unrecovered and requires human mathematical review.
\subsection{Sources}
\begin{itemize}
\item Local UnsolvedMath registry: supplied complete source record, permanent ID \(30005301\), exact source problem code \texttt{catalogue entry}, title \emph{Nonexistence of Geproci Sets in Linear General Position}. No individual proposer is identified. The supplied proposed year is 2022.
\item Supplied report citation: \emph{Mini-Workshop: Subvarieties in Projective Spaces and Their Projections} (2023), \emph{Oberwolfach Reports}, pp.~3093--3119, DOI \texttt{10.4171/owr/2022/54}: \url{https://doi.org/10.4171/owr/2022/54}. The extracted text contains the page marker 3107. The citation's 2023 date differs from the supplied proposed year and 2022 report identifier; this discrepancy is unresolved here.
\item Supplied literature reference: Chiantini et al., \emph{Configurations of points and projections}: \url{https://arxiv.org/abs/2209.04820}.
\item Supplied literature reference: Chiantini et al., \emph{Special projection properties in P4}: \url{https://arxiv.org/abs/2407.01744}.
\item Supplied literature reference: EMS Press, OWR 2022/54: \url{https://doi.org/10.4171/owr/2022/54}.
\item Literature assessment dated exactly \texttt{2026-08-22}: the supplied curation labels the problem open and states that nontrivial geproci examples in linear general position beyond four points remain an explicit open problem, while existing classifications exclude several low-degree families without proving universal nonexistence. This is supplied dated curation, not an independently verified literature assessment.
\item Supplied statement-verification note: the official report is said to confirm a narrower duplicate of the classification questions retained in \texttt{OWR-11695864-001}; no separate clean statement is promoted, and the original record remains visible for provenance. This source-verification claim is supplied registry information and has not been independently checked here.
\end{itemize}
\end{document}
